\documentclass[11pt,a4paper]{article}

% Packages
\usepackage[utf8]{inputenc}
\usepackage[T1]{fontenc}
\usepackage{lmodern}
\usepackage{graphicx}
\usepackage{booktabs}
\usepackage{array}
\usepackage{cite}
\usepackage[margin=1in]{geometry}
\usepackage{hyperref}

% =============================
% Title
% =============================
\title{
    Abnormal Human Activity Recognition in Video Surveillance: A Review
}
\author{Sanele Hlabisa \\
University of KwaZulu-Natal \\
\texttt{218023135@stu.ukzn.ac.za}}
\date{}

\begin{document}

    \maketitle

    % =============================
    % Abstract
    % =============================
    \begin{abstract}
        % Summarise the motivation, review scope, method, main themes, key gaps,
        % and overall conclusion. Write this after completing the paper.
    \end{abstract}

    % Add four to six searchable terms after completing the abstract.
    % \noindent\textbf{Keywords:} keyword one, keyword two, keyword three

    % =============================
    % Introduction
    % =============================
    \section{Introduction}
    \label{sec:introduction}
    % Introduce the review problem and explain why a synthesis of the field is needed.

    \subsection{Background and Motivation}
    \label{subsec:background-motivation}
    % Establish the importance of automated abnormal-activity recognition in
    % surveillance video and the practical motivation for reviewing the field.

    \subsection{Review Aim and Questions}
    \label{subsec:review-aim-questions}
    % State one clear aim and a small set of questions about methods, datasets,
    % evaluation, limitations, and unresolved problems.

    \subsection{Scope and Boundaries}
    \label{subsec:scope-boundaries}
    % Define what the review includes and excludes, such as video surveillance,
    % activity classification, violence detection, and video anomaly detection.

    \subsection{Paper Organization}
    \label{subsec:paper-organization}
    % Give a short roadmap of the remaining sections after the structure is stable.

    % =============================
    % Review Method
    % =============================
    \section{Review Method}
    \label{sec:review-method}
    % Explain the transparent process used to find, select, and analyse studies.

    \subsection{Search Strategy}
    \label{subsec:search-strategy}
    % Record databases, exact search terms, publication-date coverage, and the
    % date on which each search was last performed.

    \subsection{Inclusion and Exclusion Criteria}
    \label{subsec:selection-criteria}
    % Define the study types, surveillance setting, input modality, language,
    % publication type, and other criteria used to retain or reject papers.

    \subsection{Study Selection}
    \label{subsec:study-selection}
    % Describe duplicate removal, title and abstract screening, full-text
    % assessment, and the number of papers retained at each stage.

    \subsection{Data Extraction and Categorization}
    \label{subsec:data-extraction}
    % List the information collected from each paper and explain how studies are
    % grouped by task, method family, dataset, supervision, metrics, and limitations.

    % =============================
    % Domain Foundations
    % =============================
    \section{Abnormal Activity Recognition in Video Surveillance}
    \label{sec:domain-foundations}
    % Establish the concepts needed to compare the reviewed studies consistently.

    \subsection{Terminology and Task Formulations}
    \label{subsec:terminology-tasks}
    % Define abnormal activity recognition and distinguish activity classification,
    % violence detection, and video anomaly detection.

    \subsection{Surveillance-Video Challenges}
    \label{subsec:surveillance-challenges}
    % Synthesise recurring challenges such as limited labels, class imbalance,
    % low resolution, occlusion, viewpoint changes, and long untrimmed videos.

    % =============================
    % Datasets and Evaluation
    % =============================
    \section{Datasets and Evaluation}
    \label{sec:datasets-evaluation}
    % Explain how data and evaluation choices affect the comparability of results.

    \subsection{Benchmark Datasets}
    \label{subsec:benchmark-datasets}
    % Compare each relevant dataset by source, scale, classes, annotation level,
    % realism, availability, and known limitations.

    \subsection{Evaluation Metrics}
    \label{subsec:evaluation-metrics}
    % Define the metrics used in the reviewed studies and explain which task and
    % class-balance conditions each metric represents well.

    \subsection{Evaluation Protocols}
    \label{subsec:evaluation-protocols}
    % Compare data splits, cross-validation, cross-dataset testing, preprocessing,
    % and other protocol differences that influence reported performance.

    % =============================
    % Recognition Approaches
    % =============================
    \section{Recognition Approaches}
    \label{sec:recognition-approaches}
    % Organise the literature by method family and synthesise progress over time.

    \subsection{Traditional Methods}
    \label{subsec:traditional-methods}
    % Review handcrafted spatial and motion features with classical classifiers,
    % focusing on their role, strengths, and limitations.

    \subsection{CNN-Based Methods}
    \label{subsec:cnn-methods}
    % Review frame-based CNNs, transfer learning, and methods that aggregate
    % spatial features across frames.

    \subsection{3D CNN Methods}
    \label{subsec:3d-cnn-methods}
    % Explain how 3D convolutions model space and time, then compare their
    % performance, data requirements, and computational cost.

    \subsection{CNN--RNN and ConvLSTM Methods}
    \label{subsec:cnn-rnn-convlstm-methods}
    % Compare recurrent temporal models, including LRCN, LSTM, bidirectional
    % variants, and ConvLSTM, noting how spatial and temporal features interact.

    \subsection{Transformer-Based Methods}
    \label{subsec:transformer-methods}
    % Review attention-based video models, their ability to capture long-range
    % relationships, and their dependence on data and computational resources.

    % =============================
    % Comparative Analysis
    % =============================
    \section{Comparative Analysis}
    \label{sec:comparative-analysis}
    % Compare evidence across method families without treating incompatible
    % datasets, metrics, or evaluation protocols as directly equivalent.

    \subsection{Recognition Performance}
    \label{subsec:recognition-performance}
    % Compare results only under compatible settings and explain the factors that
    % may account for meaningful differences.

    \subsection{Computational Efficiency}
    \label{subsec:computational-efficiency}
    % Compare model size, operations, throughput, latency, and hardware reporting
    % where these details are available.

    \subsection{Generalization to Real Surveillance Conditions}
    \label{subsec:real-world-generalization}
    % Assess cross-scene and cross-dataset evidence, robustness to realistic
    % conditions, false alarms, and practical deployment constraints.

    % =============================
    % Research Gaps and Future Directions
    % =============================
    \section{Research Gaps and Future Directions}
    \label{sec:research-gaps}
    % Derive each gap from repeated evidence in the reviewed literature and link
    % it to a concrete research opportunity.

    \subsection{Dataset Limitations}
    \label{subsec:dataset-limitations}
    % Identify limitations in dataset size, diversity, realism, annotation quality,
    % class balance, and access that remain unresolved across studies.

    \subsection{Robustness and Generalization}
    \label{subsec:robustness-generalization}
    % Identify weaknesses in transfer across cameras, scenes, populations,
    % activity classes, and surveillance domains.

    \subsection{Computational Cost}
    \label{subsec:computational-cost}
    % Examine the gap between reported recognition performance and the efficiency
    % required for scalable or real-time surveillance deployment.

    \subsection{Explainability and Localization}
    \label{subsec:explainability-localization}
    % Assess limitations in explaining decisions and locating when and where an
    % abnormal event occurs.

    \subsection{Positioning of the PhD Research}
    \label{subsec:phd-positioning}
    % Connect the strongest evidence-backed gaps to the planned PhD contribution
    % without making unsupported novelty or performance claims.

    % =============================
    % Conclusion
    % =============================
    \section{Conclusion}
    \label{sec:conclusion}
    % Answer the review questions, summarise the strongest findings and gaps, and
    % close without introducing new evidence.

    % =============================
    % References
    % =============================
    \bibliographystyle{IEEEtran}
    \bibliography{references}

\end{document}
